\noindent\textit{2015 manuscript.}

\section{Mixtures of ideal gases}

\subsection{Gibbs paradox}


Kittel and Kroemer (1980) \cite{CKittelHKroemer1980} Problem 6 \textbf{Entropy of mixing} of Chapter 6: Ideal Gas,

Prof. Steven Anlage, Introduction to Thermodynamics and Statistical Mechanics, Spring 2011, University of Maryland Physics 404, Homework 8 Solutions, Question 9, \footnote{\url{http://www.physics.umd.edu/courses/Phys404/Anlage_Spring11/hw8solv5.pdf}}

Consider adding together the entropies of ideal gas A with initial volume $V$, and molecular mass $M_A$ (so it has a $n_{Q_A}$ quantum concentration), and of ideal gas B with initial volume $V$, molecular mass $M_B$, $n_{Q_B}$ quantum concentration:
\[
\begin{gathered}
  \sigma_{\text{initial}} = \sigma_A + \sigma_B = N_A \ln{ \left( \frac{n_{Q_A} V}{N_A} \right) } + N_A \left( 1 + \frac{d}{2} \right)  + N_B \ln{ \left( \frac{n_{Q_B} V}{N_B} \right) } + N_B \left( 1 + \frac{d}{2} \right) 
\end{gathered}
\]
After diffusive equilibrium is reached means that each of the gases now occupy a volume $2V$, gas A is free to roam in volume $2V$, as is gas B.  Thus, the entropy increases:
\[
\begin{gathered}
  \sigma_{\text{final}} = N_A \ln{ \left( \frac{n_{Q_A} 2V}{N_A} \right) } + N_A \left( 1 + \frac{d}{2} \right)  + N_B \ln{ \left( \frac{n_{Q_B} 2V}{N_B} \right) } + N_B \left( 1 + \frac{d}{2} \right) 
\end{gathered}
\]
and so
\[
\Delta \sigma = (N_A + N_B)\ln{2} = 2N\ln{2} \text{ (if $N_A=N_B=N$) }
\]

In general,
\[
\Delta \sigma = N_A \ln{ \left( \frac{V_A + V_B}{V_A} \right) } + N_B \ln{ \left( \frac{V_A + V_B}{V_A} \right) }
\]
\begin{correction}{mixing volumes}
The second logarithm is $\ln((V_A+V_B)/V_B)$.
The derivation is \S\ref{sec:binary-mixtures}, equation \eqref{eq:entropy-of-mixing}.
\end{correction}


For identical atoms $(A\equiv B)$, 
\[
\sigma_A = N \ln{ \left( \frac{ n_Q V}{N} \right) } + N \left( 1 + \frac{d}{2} \right) = \sigma_B
\]
and so
\[
\sigma_{\text{initial}} = \sigma_A + \sigma_B = 2N \ln{ \left( \frac{n_Q V}{N} \right) } + 2N\left( 1 + \frac{d}{2} \right) 
\]
For the final entropy after mixing, treat the system as a (whole) system of $2N$ indistinguishable particles occupying a volume of $2V$:
\[
\sigma_{\text{final}} = 2N \ln{ \left( \frac{n_Q(2V)}{2N} \right) } + 2N \left( 1 + \frac{d}{2} \right)
\]
and so $\Delta \sigma =0$


From Sec. 93. Mixtures of ideal gases of Landau and Lifshitz (1980) \cite{LLandauELifshitz1980}, additivity of thermodynamic quantities (such as energy and entropy) holds only if interaction between various parts of body is negligible. 

But for mixture of several substances, e.g. mixture of liquids, thermodynamic quantities are not equal to sums of thermodynamic quantities for individual components of the mixture. 

Mixtures of ideal gases are an exception, since interaction between molecules is by definition negligible.  

But for partial pressure of $i$th gas, $p_i = \frac{N_i\tau}{V} = \frac{N_i p}{N}$, $N = $ total number of molecules in mixture, $N_i$ number of molecules of $i$th gas.  

So
\[
p_j = \frac{ N_j \tau }{V} = \frac{ N_j}{N} p = X_j p 
\]
where $X_j$ is mole fraction or particle fraction.  




